\documentclass[10pt,a4paper]{amsart}
\usepackage[margin=19mm]{geometry}
\usepackage{amsmath,amssymb,mathtools}
\usepackage[scr=rsfs,cal=euler]{mathalfa}
\usepackage{xfrac}
\usepackage[unicode,hidelinks,bookmarks=false]{hyperref}
\newcommand{\ie}{i.e.\ }
\newcommand{\cf}{cf.\ }
\newcommand{\resp}{respectively\ }
\newcommand{\Subsection}{\subsection}
\providecommand{\nobreakdash}{\nobreak\hbox{-}}
\setlength{\emergencystretch}{2em}
\multlinegap0pt
\usepackage{amssymb}
\usepackage{dsfont}
\newtheorem{lem}{Lemma}
\def\R{\mathbb{R}}
\title{Stability of phase portrait for a gradient ODE with memory}
\author{Piotr Kalita, Piotr Zgliczyński}
\date{Source: arXiv:2406.00910v5 (CC BY 4.0)}
\begin{document}
\maketitle

\noindent\emph{Source and licence.} Stability of phase portrait for a gradient ODE with memory. Version arXiv:2406.00910v5 (CC BY 4.0), distributed by arXiv under the Creative Commons Attribution 4.0 International licence, http://creativecommons.org/licenses/by/4.0/, as recorded in the arXiv licence metadata for this exact version and shown on the abstract page https://arxiv.org/abs/2406.00910v5. Full licence notice: \texttt{licenses/arxiv-2406.00910-CC-BY-4.0.txt}.\par\smallskip

\noindent\textit{Editorial scope.} Counted unit: the article's lem lem:transport\_eps (source0.tex lines 1879--1892). The article's complete original proof of it is delivered with it (source0.tex lines 1893--2002, one \texttt{proof} environment of 110 lines). No mathematics is written, repaired or re-derived: the statement and the proof are the article's own lines, in the article's own notation, and no retained line is substituted or re-flowed.  Retained uncounted as the source-local context the proof needs: lines 1804--1810.  Nothing was omitted from any retained span, so the package carries no omission record.  No retained line contains a citation, so the wrapper carries no bibliography and no bibliography entry.  No retained line contains a figure environment, an image-inclusion command or drawing code, so no image is delivered and the package declares no asset dependency.  Editorial formatting only, in addition: an AMS \texttt{amsart} preamble that restates the article's own declarations and macros used by the retained lines, namely the environments \texttt{lem} on separate counters, together with the standard packages those lines need.  The retained environments print in this document as Lemma 1 in order of appearance, whereas the article prints the same items with the numbering of its own full document.  The complete native source of the article as distributed in the arXiv e-print for this exact version is bundled unchanged as \texttt{sources/160155/source0.tex}.
	 		 	\begin{lem}\label{lem:transport_disks} Let $\R^d=V+W$ where $V,W$ are two subspaces such that $\text{dim}\, V = n$ and $\text{dim}\, W = d-n$.  Let $g:B_1\to B_2$ be a $C^1$  disk, where $B_1\subset V$ and $B_2 \subset W$ are balls.  
	 		Let moreover $G:\R^d\to \R^d$ be a $C^1$ mapping.  Assume that $z_0\in \text{graph}(g)$ be such that $DG(z_0)$ is invertible and let $M$ be a $d\times d$  invertible matrix such that $M\cdot((0)_{d-n}\times \R^n) = T_{G(z_0)}G\circ (\textrm{Id}_{B_1}+g)$. Then for every $L>0$ there exists  $\delta_3, \delta_4>0$, balls $B(0,\delta_3)\subset  (0)_{d-n} \times \R^n$, $B(0,\delta_4)\subset \R^{d-n} \times (0)_n$ and the  disk $h:B(0,\delta_3)\to B(0,\delta_4)$ with the Lipschitz constant  $L$ such that 
	 		$$
	 		M \cdot \text{graph}(h) \subset G(\text{graph}(g))-G(z_0).
	 		$$
	 		Moreover $\delta_4\leq c(\delta_3)\delta_3$, where $c(\delta_3)\to 0$ as $\delta_3\to 0$. 
	 	\end{lem}

	 	\begin{lem} \label{lem:transport_eps} Let $\R^d=V+W$ where $V,W$ are two subspaces such that $\text{dim}\, V = n$ and $\text{dim}\, W = d-n$. Let moreover $X$ be a Banach space. Let $g:B_1\to B_2$ be a $C^1$  disk, where $B_1\subset V$ and $B_2 \subset W$ are balls.  Let moreover $g_\varepsilon=(g_{\varepsilon,X},g_{\varepsilon,W}):B_1\to B_2\times P$ be a family of $C^1$  disks given for $\varepsilon>0$ with $P\subset X$ being an open and bounded set such that $\sup_{v\in B_1}\left\|Dg_{\epsilon,X}(v)\right\|_{\mathcal{L}(V;X)} \leq E(\varepsilon)$, with $E(\varepsilon)\to 0$ as $\varepsilon\to 0$. Let moreover 
	 		$$
	 		\lim_{\varepsilon\to 0}\sup_{v\in B_1}\left(|g_{\varepsilon,W}(v)-g(v)|+\left|D g_{\varepsilon,W}(v)-D g(v)\right|\right) = 0.
	 		$$
	 		Next, let $G_\varepsilon=(F_{\varepsilon,X},G_{\varepsilon,d}):X \times \R^d\to X\times \R^d$ be a family of $C^1 $mappings and let $G:\R^d\to \R^d$ be a $C^1$ mapping such that
	 		$$
	 		\lim_{\varepsilon\to 0}\sup_{\eta\in P}\sup_{u\in B_1+B_2}\left(|G_{\varepsilon,d}(\eta,u)-G(u)|+\left|D_uG_{\varepsilon,d}(\eta,u)-DG(u)\right|\right)=0.
	 		$$ Moreover let $\|D G_{\varepsilon}(\eta,u)\|_{\mathcal{L}(X\times \R^d;X\times \R^d)}$ be bounded uniformly in $\varepsilon$ for $(\eta,u)$ in bounded sets in $X\times \R^d$. Assume that $z_0 = v_0+g(v_0) \in \text{graph}(g)$ is such that $DG(z_0)$  is invertible and let $M$ be a $d\times d$  invertible matrix such that $M\cdot((0)_{d-n}\times \R^n) = T_{G(z_0)}G\circ (\textrm{Id}_{B_1}+g)$. Pick $L>0$. There exist $\varepsilon_0>0$, $\delta_3, \delta_4>0$, with $\delta_4\leq c(\varepsilon_0)+c(\delta_3)\delta_3$ where $c(r)\to 0$ as $r\to 0$,  balls $B(0,\delta_3)\subset  (0)_{d-n} \times \R^n$, $B(0,\delta_4)\subset  \R^{d-n} \times (0)_n$ and a bounded set  $Q \subset X$ such that for every $\varepsilon\in (0,\varepsilon_0]$ there exists the  disk $h_\varepsilon=(h_{\varepsilon,W},h_{\varepsilon,X}):B(0,\delta_3)\to B(0,\delta_4)\times Q$ with the Lipschitz constant for $h_{\varepsilon,W}$ equal to $L$ and 
	 		$$
	 		 M \cdot \text{graph}(h_{\varepsilon,W}) \subset G_{\varepsilon,d}(\text{graph}(g_\varepsilon))-G(z_0).
$$
Moreover $h_{\varepsilon,X}(0,\overline{v}) = G_{\varepsilon,X}(g_{\varepsilon,X}(v),v+g_{\varepsilon,W}(v))$, where $v$ is such that $	 		 M \cdot ((0,\overline{v})+h_{\varepsilon,W}(0,\overline{v})) =  G_{\varepsilon,d}(g_{\varepsilon,X}(v),v+g_{\varepsilon,W}(v)))-G(z_0)
$.
	 		\end{lem}

	 		\begin{proof}
	 			In the lemma statement we denote $\textrm{graph}(g_\varepsilon) = \{ (g_{\varepsilon,X}(p),v+g_{\varepsilon,W}(p))\,:\ p\in B_1 \}$. Let $B\subset V$ be a ball such that $v_0+B\subset B_1$.
 If $v\in B$, then the point in the graph of $g_\varepsilon$ is denoted by
	 			\begin{align*}
	 				&  (g_{\varepsilon,X}(v+v_0),v+v_0+g_{\varepsilon,W}(v+v_0)) \\
	 				& \ \ =    \left(g_{\varepsilon,X}(v_0)+Dg_{\varepsilon,X}(v_0)v+\Delta_1(v),v_0+v+g(v_0)+g_{\varepsilon,W}(v_0)-g(v_0)+
	 				Dg_{\varepsilon,W}(v_0)v+\Delta_2(v)\right)\\
	 				& \ \ = (g_{\varepsilon,X}(v_0),z_0) + \left(Dg_{\varepsilon,X}(v_0)v,v+Dg_{\varepsilon,W}(v_0)v\right) + (0,g_{\varepsilon,W}(v_0)-g(v_0)) +  \left(\Delta_1(v),\Delta_2(v)\right)\\
	 				& \ \  = (g_{\varepsilon,X}(v_0),z_0) + \left(Dg_{\varepsilon,X}(v_0)v,v+Dg(v_0)v\right) +   \left(0,\left(Dg_{\varepsilon,W}(v_0)-Dg(v_0)\right)v\right) \\
	 				& \qquad \qquad + (0,g_{\varepsilon,W}(v_0)-g(v_0)) +  \left(\Delta_1(v),\Delta_2(v)\right)\\
	 				& \ \  = (g_{\varepsilon,X}(v_0),z_0) + I + II + III + IV.
	 			\end{align*}
	 			Now, consider the image of this point by the map $G_\varepsilon=(G_{\varepsilon,X},G_{\varepsilon,d})$.
	 			\begin{align*}
	 				& G_{\varepsilon,d}(g_{\varepsilon,X}(v+v_0),v+v_0+g_{\varepsilon,W}(v+v_0)) \\
	 				& \  = G(z_0) + (G_{\varepsilon,d}(g_{\varepsilon,X}(v_0),z_0)-G(z_0)) + D G_{\varepsilon,d}(g_{\varepsilon,X}(v_0),z_0)I \\
	 				& \quad + DG_{\varepsilon,d}(g_{\varepsilon,X}(v_0),z_0)(II+III+IV)\\
	 				& \quad +(\overline{DG_{\varepsilon,d}([(g_{\varepsilon,X}(v_0),z_0),(g_{\varepsilon,X}(v+v_0),v+v_0+g_{\varepsilon,W}(v+v_0))]}-DG_{\varepsilon,d}(g_{\varepsilon,X}(v_0),z_0))\cdot\\
	 				& \qquad \qquad \cdot(I+II+III+IV)\\
	 				& \  = G(z_0) + DG(z_0)\left(v+Dg(v_0)v\right)  + \left(D G_{\varepsilon,d}(g_{\varepsilon,X}(v_0),z_0)-DG(z_0)\right)I + \\
	 				& \quad  +  (G_{\varepsilon,d}(g_{\varepsilon,X}(v_0),z_0)-G(z_0)) + D G_{\varepsilon,d}(g_{\varepsilon,X}(v_0),z_0)(II+III+IV)\\
	 				& \quad +(\overline{DG_{\varepsilon,d}([(g_{\varepsilon,X}(v_0),z_0),(g_{\varepsilon,X}(v+v_0),v+v_0+g_{\varepsilon,W}(v+v_0))]}-DG_{\varepsilon,d}(g_{\varepsilon,X}(v_0),z_0))\cdot\\
	 				& \qquad \qquad \cdot(I+II+III+IV).
	 			\end{align*}
	 			We can rewrite the above equation as 
	 			$$
	 			G_{\varepsilon,d}(g_{\varepsilon,X}(v+v_0),v+v_0+g_{\varepsilon,W}(v+v_0)) - G(z_0) = DG(z_0)\left(v+Dg(v_0)v\right)  + \Delta_3 + \Delta_4+\Delta_5 + \Delta_6. 
	 			$$
	 			As in Lemma \ref{lem:transport_disks}  we must prove that for every ball $B(0,\delta)\subset B$ there exists $\delta_3(\delta)$ such that for every $(0,\overline{v}) \in B(0,\delta_3)$ we can find $v\in B(0,\delta)$ such that
	 			$$
	 			M\cdot(0,\overline{v}) = DG(z_0)\left(v+Dg(v_0)v\right)  + \Pi_{M\cdot(0,\overline{v})}(\Delta_3 + \Delta_4+\Delta_5 + \Delta_6).
	 			$$
	 			The proof follows, again, by the homotopy invariance of the Brouwer degree as in the previous Lemma. Indeed, there exist constants $c(\varepsilon)\to 0$ as $\varepsilon\to 0$ and $c(|v|)\to 0$ as $|v|\to 0$, for which we obtain   
	 			\begin{align*}	 				
	 			& |\Delta_3| = |\left(D G_{\varepsilon,d}(g_{\varepsilon,X}(v_0),z_0)-DG(z_0)\right)I| \leq c(\varepsilon) |I| \leq c(\varepsilon)|v|, \\
	 			& |\Delta_4| = |(G_{\varepsilon,d}(g_{\varepsilon,X}(v_0),z_0)-G(z_0))| \leq c(\varepsilon),\\
	 			& |\Delta_5| = |D G_{\varepsilon,d}(g_{\varepsilon,X}(v_0),z_0)(II+III+IV)| \leq C|II+III+IV| \leq c(\varepsilon)+(c(\varepsilon)+c(|v|))|v|,\\
	 			& |\Delta_6| \leq c(|v|)(|I+II+III+IV|)\leq c(|v|)( |v| + c(\varepsilon)).
	 			\end{align*}
	 			Thus 
	 			$$
	 			|\Delta_3+\Delta_4+\Delta_5+\Delta_6| \leq c(\varepsilon) + c(|v|)|v|,
	 			$$ 
	 			 and the homotopy to be used is. 
	 			$$
	 			B(0,\delta)\times [0,1]\ni (v,\theta) \mapsto M^{-1}(DG(z_0)\left(v+Dg(v_0)v\right)  + \theta \Pi_{M\cdot(0,\overline{v})}(\Delta_3 + \Delta_4+\Delta_5 + \Delta_6)).
	 			$$
	 			Now, consider two points in the graph of $g_\varepsilon$. Denote them by
	 			$(\eta_1,z_1) = (g_{\varepsilon,X}(v_1+v_0),v_1+v_0 + g_{\varepsilon,W}(v_1+v_0))$ and $(\eta_2,z_2) = (g_{\varepsilon,X}(v_2+v_0),v_2+v_0 + g_{\varepsilon,W}(v_2+v_0))$. Now
	 			$
	 			G_{\varepsilon,d}(\eta_1,z_1) = M\cdot (\overline{w}_1,\overline{v}_1)
	 			$ and $
	 			G_{\varepsilon,d}(\eta_2,z_2) = M\cdot (\overline{w}_2,\overline{v}_2).
	 			$
	 			The function $G_{\varepsilon,d}:X\times \R^d \to \R^d$ will be treated as the function $G_{\varepsilon,d}:X\times V\times W \to \R^d$ and the corresponding derivatives will be denoted by $D_\eta G_{\varepsilon,d}$, $D_v G_{\varepsilon,d}$, and $D_w G_{\varepsilon,d}$. We calculate
	 			\begin{align*}
	 				& M\cdot (\overline{w}_2-\overline{w}_1,\overline{v}_2-\overline{v}_1) =  G_{\varepsilon,d}(\eta_2,z_2) - G_{\varepsilon,d}(\eta_1,z_1) \\
	 				& \ \in
	 				\overline{D_v G_{\varepsilon,d}([(\eta_1,z_1),(\eta_2,z_2)])}(v_2-v_1)\\
	 				& \qquad \qquad + \overline{D_w G_{\varepsilon,d}([(\eta_1,z_1),(\eta_2,z_2)])}(g_{\varepsilon,W}(v_2+v_0)-g_{\varepsilon,W}(v_1+v_0))\\
	 				& \qquad\qquad  + \overline{D_\eta G_{\varepsilon,d}([(\eta_1,z_1),(\eta_2,z_2)])}(g_{\varepsilon,X}(v_2+v_0)-g_{\varepsilon,X}(v_1+v_0))\\
	 				& \ \subset
\overline{D_v G_{\varepsilon,d}([(\eta_1,z_1),(\eta_2,z_2)])}(v_2-v_1) + \overline{D_w G_{\varepsilon,d}([(\eta_1,z_1),(\eta_2,z_2)])}\ \overline{D g_{\varepsilon,W}([v_2+v_0,v_1+v_0])} (v_2-v_1)\\
	 				& \qquad \qquad  + \overline{D_\eta G_{\varepsilon,d}([(\eta_1,z_1),(\eta_2,z_2)])}\ \overline{D g_{\varepsilon,X}([v_2+v_0,v_1+v_0])} (v_2-v_1).
	 			\end{align*}
	 			Furthermore, treating $G:\R^d\times \R^d$ as the function $G:X\times \R^d\times \R^d$, independent on the variable in $X$, we obtain  
	 			\begin{align*}
	 				& M\cdot (\overline{w}_2-\overline{w}_1,\overline{v}_2-\overline{v}_1) \in D_v G(z_0)\left(Dg(v_0)(v_2-v_1)+v_2-v_1\right)\\
	 				& \qquad +
	 				\overline{(D_vG_{\varepsilon,d}(\cdot)-D_vG(\cdot))([(\eta_1,z_1),(\eta_2,z_2)])}(v_2-v_1) + \overline{D_vG([z_1,z_2])-D_vG(z_0)}(v_2-v_1)\\
	 				& \qquad + \overline{(D_w G_{\varepsilon,d}(\cdot)-D_w G (\cdot))([(\eta_1,z_1),(\eta_2,z_2)])}\ \overline{D g_{\varepsilon,W}([v_2+v_0,v_1+v_0])}(v_2-v_1)\\
	 				& \qquad + \left(\overline{D_w G ([z_1,z_2])}-D_wG(z_0)\right)\overline{D g_{\varepsilon,W}([v_2+v_0,v_1+v_0])}(v_2-v_1)\\
	 				& \qquad + D_wG(z_0)\overline{(D g_{\varepsilon,W}(\cdot)-Dg(\cdot))([v_2+v_0,v_1+v_0])}(v_2-v_1)\\
	 				& \qquad +D_wG(z_0)\left(\overline{Dg([v_2+v_0,v_1+v_0])}-Dg(v_0)\right)(v_2-v_1)\\
	 				& \qquad  + \overline{D_\eta G_{\varepsilon,d}([(\eta_1,z_1),(\eta_2,z_2)])}\ \overline{D g_{\varepsilon,X}([v_2+v_0,v_1+v_0])} (v_2-v_1) \\
	 				& \ \ = I + II+III+IV+V+VI+VII+VIII.
	 			\end{align*}
	 			The  term $I$ in the last sum belongs to the tangent space $T_{G(z_0)}G\circ(Id_{B_1}+g)$. Terms $II, IV, VI$ satisfy the estimate $|II+IV+VI|\leq C_2(\varepsilon)|v_2-v_1|$ with $C_2(\varepsilon)\to 0$ as $\varepsilon\to 0$. Terms $III,V,VII$ satisfy the bound $|III+V+VII|\leq C_3(\delta_3)|v_2-v_1|$, where $C_3(\delta_3)$ also tends to zero as $\delta_3\to 0$. As for the term $VIII$ we have the bound $|VIII|\leq C_4 E(\varepsilon)|v_2-v_1|$, where $C_4$ is a bound on $\|D_\eta G_{\varepsilon,d}(\eta,u)\|_{\mathcal{L}(X;\R^d)}$ for $\eta,u\in P\times (B_1+B_2)$.
	 			It follows that
	 			$$
	 			|(0,\overline{v}_2-\overline{v}_1) | \geq C_1 |v_2-v_1| - (C_2(\varepsilon)+C_3(\delta_3)+C_4E(\varepsilon))|v_2-v_1|,
	 			$$
	 			and
	 			$$
	 			|(\overline{w}_2-\overline{w}_1,0) | \leq  (C_2(\varepsilon)+C_3(\delta_3)+C_4E(\varepsilon))|v_2-v_1|.
	 			$$
 Hence
	 			$$
	 			|(\overline{w}_2-\overline{w}_1,0) | \leq  \frac{C_2(\varepsilon)+C_3(\delta_3)+C_4E(\varepsilon)}{C_1-C_2(\varepsilon)-C_3(\delta_3)-C_4E(\varepsilon)}|(0,\overline{v}_2-\overline{v}_1) |,
	 			$$
	 			and we have the required Lipschitz condition with the arbitrarily small constant obtained by decreasing $\varepsilon$ and the radius $\delta_3$.
	 			Finally, let us estimate
	 			\begin{align*}
	 				& G_{\varepsilon,X}(\eta_2,z_2) - G_{\varepsilon,X}(\eta_1,z_1) \\
	 				& \ \in
	 				\overline{D_v G_{\varepsilon,X} ([(\eta_2,z_2),(\eta_1,z_1)])}(v_2-v_1)\\
	 				& \qquad \qquad + \overline{D_w G_{\varepsilon,X} ([(\eta_2,z_2),(\eta_1,z_1)])}(g_{\varepsilon,W}(v_2+v_0)-g_{\varepsilon,W}(v_1+v_0)) \\
	 				& \qquad \qquad  + \overline{D_\eta G_{\varepsilon,X} ([(\eta_2,z_2),(\eta_1,z_1)])}(g_{\varepsilon,X}(v_2+v_0)-g_{\varepsilon,X}(v_1+v_0))\\
	 				& \ \subset 
	 								\overline{D_v G_{\varepsilon,X} ([(\eta_2,z_2),(\eta_1,z_1)])}(v_2-v_1)\\
	 								& \qquad \qquad + \overline{D_w G_{\varepsilon,X} ([(\eta_2,z_2),(\eta_1,z_1)])}\ \overline{D g_{\varepsilon,W}([v_2+v_0,v_1+v_0])}(v_2-v_1)\\
	 				& \qquad \qquad  + \overline{D_\eta G_{\varepsilon,X} ([(\eta_2,z_2),(\eta_1,z_1)])}\ \overline{D g_{\varepsilon,X}([v_2+v_0,v_1+v_0])}(x_2-x_1).
	 			\end{align*}
	 			This means that
	 			$$
	 			\|G_{\varepsilon,X}(\eta_2,z_2) - G_{\varepsilon,X}(\eta_1,z_1)\|_X \leq L_1|v_2-v_1| \leq  L_2|(0,\overline{v}_2-\overline{v}_1) |,
	 			$$
	 			which completes the proof.
	 			
	 			\end{proof}

\end{document}
